\section{Related Work}
\label{sec:related_work}

\subsection{Air Quality Monitoring with IoT Sensor Networks}

The deployment of low-cost sensor networks for air quality monitoring has accelerated in recent years, driven by advances in microelectromechanical systems (MEMS) and IoT connectivity. Karmakar et al.~(2023) developed the DALTON platform and conducted a comprehensive field study across 28 indoor sites, revealing that conventional indices such as the Indoor Air Quality Index (IAQI) fail to capture complex spatio-temporal pollution dynamics driven by occupant activities such as cooking, cleaning, and ventilation changes \cite{karmakar2024exploring}. Their proposed Healthy Home Index (HHI) integrates average pollutant exposure, statistical variation, and spatial spread over a 10-minute sliding window. R\'{o}denas Garc\'{i}a et al.~(2022) catalogued the specifications and limitations of commercially available low-cost sensors, highlighting trade-offs in detection range, measurement tolerance, response time, and power consumption \cite{rodenas2022review}. More recently, innovations in IoT-enabled monitoring systems have incorporated machine learning calibration and multi-sensor fusion for improved accuracy \cite{innovations2025sensors}, while TinyML-based approaches have demonstrated real-time ozone prediction on Arduino Nano microcontrollers with $R^2$ values of 0.95 \cite{akhyar2025tinyml}.

\subsection{Sensor Data Reduction and Compression}

Reducing the volume of sensor data prior to transmission is a well-established strategy for conserving energy in wireless sensor networks. Atitallah et al.~(2023) proposed an autonomous IoT data reduction method based on adaptive thresholds, incorporating concept drift detection to dynamically balance reduction rate with reconstruction accuracy \cite{atitallah2023autonomous}. Their approach achieved an 11.7\% improvement in accuracy for the same reduction rate, or equivalently a 17.3\% improvement in reduction rate for the same accuracy, when compared with fixed-threshold baselines such as Kalman filtering, LMS, and PIP algorithms. SAX-LZW two-level compression strategies have demonstrated significant data reduction for IoT sensor streams by first discretising continuous signals into symbolic representations and then applying dictionary-based compression \cite{aljanabi2022sax}. Bounded-error pruning approaches selectively discard sensor readings that fall within a configurable tolerance band, achieving energy savings proportional to the signal's stationarity \cite{azar2020bounded}. However, these methods operate at the raw signal level without considering the downstream requirements of inference models deployed on the same edge device.

\subsection{Neural Network Quantisation for Edge Deployment}

Quantisation reduces the numerical precision of neural network weights and activations from 32-bit floating point to lower bit-widths, yielding proportional reductions in memory footprint, computational cost, and energy consumption. Nagel et al.~(2021) provided a comprehensive white paper on neural network quantisation, establishing theoretical foundations for uniform and non-uniform quantisation, calibration techniques, and the trade-off between bit-width and task accuracy \cite{nagel2021whitepaper}. Post-training quantisation (PTQ) methods, which require no retraining, have shown significant progress: ZeroQuant~(2022) introduced token-wise quantisation for transformers \cite{zeroquant2022}; SmoothQuant~(2024) demonstrated training-free W8A8 quantisation through activation-weight smoothing \cite{smoothquant2024}; and QDrop~(2023) randomly dropped quantisation operations during calibration to improve extremely low-bit PTQ \cite{qdrop2023}.

Quantisation-aware training (QAT) methods incorporate quantisation effects into the training loop for superior accuracy at very low bit-widths. EdgeQAT~(2024) introduced entropy and distribution guided training specifically designed for lightweight large language models on edge devices \cite{edgeqat2024}. EfficientQAT~(2024) achieved state-of-the-art results by quantising weights and activations to 4 bits with minimal accuracy loss \cite{efficientqat2024}. For mixed-precision inference, DyBit~(2023) proposed dynamic bit-precision numbers that adapt precision to the statistical properties of weights and activations, achieving 1.997\% higher accuracy than prior methods \cite{dybit2023}. HADQ-Net~(2025) presented a hardware-adaptive QAT framework demonstrating that quantisation from FP32 to lower bit-widths significantly reduces memory footprint and energy consumption while maintaining accuracy across classification, detection, and segmentation tasks \cite{hadqnet2025}.

On edge hardware, Palakodety et al.~(2025) achieved a 14.5$\times$ speedup with INT8 PTQ on the Sipeed MaixCam platform (RISC-V, 1~TOPS NPU) while maintaining high mAP at only 1.96~W average power \cite{palakodety2025tpu}. Guo et al.~(2020) demonstrated that 4-bit fixed-point LeNet-5 achieves 99.28\% accuracy with an energy efficiency of 213~GOP/s/W on FPGA \cite{guo2020fixedpoint}. The LUTIN framework~(2024) converted matrix multiplication to table lookups, achieving 2.34$\times$ speedup and 2.04$\times$ energy efficiency improvement by investigating bit-widths below 8 bits \cite{lutin2024}.

\subsection{ADC Energy Modelling}

The relationship between ADC resolution and energy consumption is well characterised by the Walden figure of merit (FoM), defined as:
\begin{equation}
    \text{FoM}_W = \frac{P_{\text{ADC}}}{2^{\text{ENOB}} \cdot f_s}
    \label{eq:walden}
\end{equation}
where $P_{\text{ADC}}$ is the ADC power consumption, ENOB is the effective number of bits, and $f_s$ is the sampling frequency \cite{andrulis2024adc}. This implies that power scales exponentially with resolution: $P_{\text{ADC}} = \text{FoM}_W \cdot 2^N \cdot f_s$. Andrulis et al.~(2024) developed an open-source architecture-level ADC energy and area model showing that ADC energy depends on throughput, ENOB, and technology node, with energy increasing exponentially with ENOB \cite{andrulis2024adc}. State-of-the-art SAR ADCs achieve FoMs ranging from 0.44~fJ/conv-step (11-bit, 90~nm CMOS) \cite{waldenADC2018} to 16.6~fJ/conv-step (8-bit, 5~nm CMOS). For a reconfigurable 2--8 bit ADC designed for computing-in-memory, adjustable resolution allows unnecessary power consumption to be eliminated during low-resolution requirements \cite{reconfigADC2025}.

\subsection{Research Gap}

The existing literature treats sensor-level data reduction and inference-level model quantisation as disjoint problems. No prior work has proposed a unified framework that jointly optimises ADC bit-width allocation and neural network quantisation for air quality monitoring on edge devices. Furthermore, existing adaptive sampling methods use static or slowly adapting thresholds that do not exploit the strong periodic structure of indoor pollutant data---characterised by long stable periods (nighttime, unoccupied periods) interrupted by sharp, activity-driven transients (cooking, cleaning, ventilation changes) \cite{karmakar2024exploring}. The VABQ framework addresses this gap by introducing a variance-entropy criterion that governs both sensing and inference quantisation within a single mathematically grounded pipeline.
